% =====================================================================
% preamble.tex
% 主文档预导言：宏包加载、字体配置、全局设置
%
% 编译：xelatex （必须 xelatex/lualatex，不能 pdflatex —— 因为 ctex+xeCJK）
%
% Wave 0 - 不允许后续 Wave subagents 修改本文件
% =====================================================================

% --- 字体与中文 ---
\usepackage[UTF8, scheme=plain, fontset=none]{ctex}
\setCJKmainfont[
  BoldFont = Noto Serif CJK SC Bold,
  ItalicFont = Noto Serif CJK SC,
]{Noto Serif CJK SC}
\setCJKsansfont{Noto Sans CJK SC}
\setCJKmonofont{Noto Sans Mono CJK SC}
\setmainfont{TeX Gyre Termes}      % 西文衬线
\setsansfont{TeX Gyre Heros}       % 西文无衬线
\setmonofont{DejaVu Sans Mono}[Scale=0.92]

% --- 几何 ---
\usepackage[a4paper, margin=2.5cm, headheight=14pt]{geometry}
\usepackage{fancyhdr}
\pagestyle{fancy}
\fancyhf{}
\fancyhead[LE,RO]{\thepage}
\fancyhead[LO]{\nouppercase\rightmark}
\fancyhead[RE]{\nouppercase\leftmark}
\renewcommand{\headrulewidth}{0.4pt}

% --- 数学 ---
\usepackage{amsmath, amssymb, amsthm, mathtools}
\usepackage{bm}
\usepackage{physics}      % \bra, \ket, \dv, \pdv, etc.
\usepackage{empheq}       % grouped equations (used by some chapters)
\usepackage{enumitem}     % enumerate options like [leftmargin=...]
\usepackage{siunitx}

% --- 图形与表格 ---
\usepackage{graphicx}
\usepackage{float}        % H float specifier
\usepackage{tikz}
\usetikzlibrary{arrows.meta, calc, positioning, shapes.geometric, fit, backgrounds, decorations.pathreplacing, trees, matrix}
\usepackage{pgfplots}
\pgfplotsset{compat=1.18}
\usepackage{booktabs}
\usepackage{array}
\usepackage{multirow}
\usepackage{tabularx}
\usepackage{longtable}

% --- 代码 listings ---
\usepackage{listings}
\usepackage[dvipsnames]{xcolor}
% =====================================================================
% code_macros.tex
% listings 包风格定义：cpp / fortran / python
%
% 用法：subagents 通过 \lstinputlisting[style=cpp]{code_excerpts/snippets/foo.cpp}
%       引用 extract.py 抽取的真实代码片段
% =====================================================================

\definecolor{lstbg}{RGB}{248,248,248}
\definecolor{lstkw}{RGB}{0,0,180}
\definecolor{lstcomment}{RGB}{60,120,60}
\definecolor{lststring}{RGB}{160,80,0}
\definecolor{lstnum}{RGB}{120,120,120}
\definecolor{lstrule}{RGB}{200,200,200}

\lstset{
  basicstyle=\ttfamily\footnotesize,
  backgroundcolor=\color{lstbg},
  keywordstyle=\color{lstkw}\bfseries,
  commentstyle=\color{lstcomment}\itshape,
  stringstyle=\color{lststring},
  numbers=left,
  numberstyle=\tiny\color{lstnum},
  numbersep=8pt,
  frame=single,
  framerule=0.4pt,
  rulecolor=\color{lstrule},
  breaklines=true,
  breakatwhitespace=false,
  tabsize=4,
  showstringspaces=false,
  captionpos=b,
  keepspaces=true,
  columns=fullflexible,
}

% C++ 风格
\lstdefinestyle{cpp}{
  language=C++,
  morekeywords={uint64_t, int64_t, size_t, auto, constexpr, nullptr, override, final, noexcept},
}

% Fortran 90 风格
\lstdefinestyle{f90}{
  language=[90]Fortran,
  morekeywords={module, contains, implicit, intent, allocatable, pure, elemental, subroutine, end},
}

% Fortran 77 风格
\lstdefinestyle{f77}{
  language=[77]Fortran,
}

% Python 风格
\lstdefinestyle{python}{
  language=Python,
  morekeywords={async, await, match, case},
}

% YAML / 输入文件风格（plain，关键字按需）
\lstdefinestyle{config}{
  language=,
  morekeywords={true, false},
  showstringspaces=false,
}


% --- 算法伪码 ---
\usepackage[ruled, vlined, linesnumbered]{algorithm2e}

% --- 引用与超链接 ---
\usepackage[colorlinks=true, linkcolor=NavyBlue, citecolor=ForestGreen, urlcolor=NavyBlue]{hyperref}
\usepackage[capitalise, noabbrev]{cleveref}
\crefname{equation}{式}{式}
\crefname{figure}{图}{图}
\crefname{table}{表}{表}
\crefname{section}{节}{节}
\crefname{chapter}{章}{章}
\crefname{algorithm}{算法}{算法}
\crefname{listing}{清单}{清单}

% --- 参考文献 ---
\usepackage[backend=biber, style=numeric-comp, sorting=none]{biblatex}
\addbibresource{references.bib}

% --- 定理环境 ---
\theoremstyle{definition}
\newtheorem{definition}{定义}[chapter]
\newtheorem{example}{例}[chapter]
\newtheorem{remark}{注}[chapter]

\theoremstyle{plain}
\newtheorem{theorem}{定理}[chapter]
\newtheorem{proposition}{命题}[chapter]
\newtheorem{lemma}{引理}[chapter]
\newtheorem{corollary}{推论}[chapter]

% --- 自定义环境 ---
\usepackage{tcolorbox}
\tcbuselibrary{breakable, skins}

% 数字闭环框（§x.6 用）
\newtcolorbox{numericalbox}[1][]{
  enhanced, breakable,
  colback=blue!3, colframe=blue!50!black, boxrule=0.6pt,
  title={\textbf{数字闭环}}, fonttitle=\bfseries,
  #1
}

% 陷阱框（§x.8 用）
\newtcolorbox{pitfallbox}[1][]{
  enhanced, breakable,
  colback=red!3, colframe=red!60!black, boxrule=0.6pt,
  title={\textbf{常见陷阱}}, fonttitle=\bfseries,
  #1
}

% 代码落地框（§x.5 用，可选）
\newtcolorbox{codelandbox}[1][]{
  enhanced, breakable,
  colback=gray!3, colframe=black!60, boxrule=0.4pt,
  title={\textbf{代码落地}}, fonttitle=\bfseries,
  #1
}

% --- 物理符号宏 ---
% =====================================================================
% physics_macros.tex
% 物理符号宏定义
%
% 维护规则：
%   - 跨章节复用的符号必须在此定义，确保全书一致性
%   - Wave 1+ subagents 只能 *添加* 新宏，不能修改/删除现有宏
%   - 添加新宏时在 commit message 注明
% =====================================================================

% --- 基本物理量 ---
\newcommand{\hbarc}{\hbar c}                            % ℏc
\newcommand{\Tlab}{T_{\text{lab}}}                      % 实验室能量
\newcommand{\Tcm}{T_{\text{cm}}}                        % 质心系能量
\newcommand{\Mn}{M_n}                                   % 中子质量
\newcommand{\Mp}{M_p}                                   % 质子质量
\newcommand{\Md}{M_d}                                   % 氘核质量
\newcommand{\Ed}{E_d}                                   % 氘核结合能

% --- 网格 / 离散化 ---
\newcommand{\Np}{N_p}                                   % p 网格点数
\newcommand{\Nq}{N_q}                                   % q 网格点数
\newcommand{\NpWP}{N_{p}^{\text{WP}}}                   % WP bin 数
\newcommand{\NqWP}{N_{q}^{\text{WP}}}
\newcommand{\Nalpha}{N_\alpha}                          % 部分波通道数

% --- 算符 ---
\newcommand{\Top}{\hat{T}}                              % T-算符
\newcommand{\Vop}{\hat{V}}                              % 相互作用
\newcommand{\Hop}{\hat{H}}                              % 哈密顿量
\newcommand{\Hzero}{\hat{H}_0}                          % 自由哈密顿量
\newcommand{\Gop}{\hat{G}}                              % Green 算符
\newcommand{\Gzero}{\hat{G}_0}                          % 自由 Green 算符
\newcommand{\Pop}{\hat{\mathcal{P}}}                    % 置换 P123
\newcommand{\Uop}{\hat{U}}                              % AGS U-算符
\newcommand{\Aop}{\hat{\mathcal{A}}}                    % 反对称化算符

% --- Jacobi 动量 / 能量 ---
\newcommand{\jacobip}{p}
\newcommand{\jacobiq}{q}
\newcommand{\bvec}[1]{\bm{#1}}
\newcommand{\jacvec}{\bvec{p}, \bvec{q}}
\newcommand{\Ejacobi}{E_{pq}}

% --- Jacobi 态与部分波态 ---
\newcommand{\jacobistate}[2]{\ket{#1, #2}}              % |p, q⟩
\newcommand{\pwstate}[1]{\ket{#1}}                      % |α⟩
\newcommand{\pwdual}[1]{\bra{#1}}                       % ⟨α|
\newcommand{\WPbin}[1]{\ket{x_{#1}}}                    % WP 第 i 个 bin

% --- 部分波索引 ---
\newcommand{\PWchan}{\alpha}
\newcommand{\PWchanp}{\alpha'}
\newcommand{\PWtuple}{(\ell, s, j; \lambda, I; J^\pi, T)}

% --- AGS / Faddeev ---
\newcommand{\Uampl}[2]{U_{#1#2}}                        % U_{αβ}
\newcommand{\Tampl}[2]{T_{#1#2}}                        % T_{αβ}

% --- 散射观测量 ---
\newcommand{\dsigma}{d\sigma}
\newcommand{\dOmega}{d\Omega}
\newcommand{\diffXS}{\frac{\dsigma}{\dOmega}}           % 微分截面
\newcommand{\iTeleven}{i T_{11}}
\newcommand{\Ttwentyzero}{T_{20}}
\newcommand{\Ttwentyone}{T_{21}}
\newcommand{\Ttwentytwo}{T_{22}}
\newcommand{\Ay}{A_y}                                   % 极化分析能力

% --- 3NF 低能常数 ---
\newcommand{\cD}{c_D}
\newcommand{\cE}{c_E}
\newcommand{\cone}{c_1}
\newcommand{\cthree}{c_3}
\newcommand{\cfour}{c_4}

% --- 数学符号 ---
\newcommand{\re}{\operatorname{Re}}
\newcommand{\im}{\operatorname{Im}}
\newcommand{\sgn}{\operatorname{sgn}}
\newcommand{\eps}{\varepsilon}
\newcommand{\nudif}{n_{\text{dif}}}                      % Padé 阶差
\newcommand{\TF}{\mathcal{F}}                            % Fourier 变换
\newcommand{\identity}{\mathbb{1}}


% --- 章节深度 ---
\setcounter{tocdepth}{2}
\setcounter{secnumdepth}{3}
